\documentclass[11pt]{article}
\usepackage[T1]{fontenc}
\usepackage[margin=1in]{geometry}
\usepackage{amsmath,amssymb,amsthm,mathtools}
\usepackage[colorlinks=true,linkcolor=blue,citecolor=blue,urlcolor=blue]{hyperref}
\newcommand{\ddc}{dd^c}
\newcommand{\PSH}{\operatorname{PSH}}
\newcommand{\tr}{\operatorname{tr}}
\newcommand{\cA}{\mathcal A}
\newcommand{\dV}{\mathrm dV}
\newtheorem{theorem}{Theorem}
\newtheorem{lemma}[theorem]{Lemma}
\newtheorem{proposition}[theorem]{Proposition}
\newtheorem{corollary}[theorem]{Corollary}
\theoremstyle{remark}
\newtheorem{remark}[theorem]{Remark}
\title{Supporting probabilities for a continuous positive complex Hessian}
\author{Research programme on non-K\"ahler pluripotential theory}
\date{8 October 2026: unverified restricted draft}
\begin{document}
\maketitle
\begin{abstract}
We give a sufficient condition for uniqueness of the supporting probability
of the Hermitian Monge--Amp\`ere normalization scalar: an actual bounded
solution has a continuous, everywhere strictly positive complex Hessian.
The proof linearizes along bounded admissible potentials and establishes
the required Fredholm statement for a continuous uniformly elliptic
operator on $W^{2,p}$. It needs no H\"older or Dini modulus for the Hessian.
Product examples give the supporting probability explicitly for every
positive continuous density depending on the curve factor, including a
non-Dini example on a non-K\"ahler threefold. The unconditional
every-continuous-density question in P16 remains open.
\end{abstract}

\paragraph{Status and scope.}
This author-developed draft has not passed independent reconstruction or
final mathematical review. The elliptic arguments use classical analysis;
no claim of literature priority is made. The product-of-two-curves
specialization overlaps the nonsingular specialization of the previously
checked product-adjoint argument \cite{S07}. The general continuous-Hessian
criterion, rather than that specialization, is the proposed restricted
result. Supporting-set nonemptiness and its $L^q$ bounds are imported
background of the original problem.

\section{Statement and conventions}
Let $X$ be a connected compact complex manifold of dimension $n\geq2$,
$\omega$ a smooth positive real $(1,1)$-form and $\dV$ a smooth positive
probability volume. Set $\ddc=i\partial\bar\partial/(2\pi)$.
Every quasi-plurisubharmonic function denotes its genuine canonical upper
representative. For $F\in C(X,\mathbb R)$ let $c_0(F)>0$ be the unique
normalization scalar for an actual bounded solution of
\begin{equation}\label{eq:equation}
 (\omega+\ddc u)^n=c_0(F)e^F\dV,
 \qquad u\in\PSH(X,\omega)\cap L^\infty(X).
\end{equation}
Bounded products are the local Bedford--Taylor products with the finite
smooth-background expansion. Put $B_0(F)=\log c_0(F)$ and
\begin{equation}\label{eq:support}
 \cA_0(F)=\left\{\nu\text{ a Borel probability}:
 B_0(G)\leq B_0(F)-\int_X(G-F)\,\mathrm d\nu
 \text{ for all }G\in C(X,\mathbb R)\right\}.
\end{equation}
We use the problem's known background: $B_0$ is concave, decreasing and
obeys $B_0(F+a)=B_0(F)-a$; $\cA_0(F)$ is nonempty, and its members have
$L^q(\dV)$ densities for every $1<q<n/(n-1)$. These assertions are not
new conclusions of this article.

\begin{theorem}\label{thm:criterion}
Suppose a bounded genuine solution $u$ of \eqref{eq:equation} has
distributional complex Hessian represented by a continuous real form
\[
 \sigma=\omega+\ddc u>0
\]
everywhere on $X$. Then $\cA_0(F)$ consists of one probability $\nu_F$.
For every $H\in C(X,\mathbb R)$,
\[
 \lim_{t\to0}\frac{B_0(F+tH)-B_0(F)}{t}
       =-\int_XH\,\mathrm d\nu_F.
\]
Neither a full real $C^2$ representative of $u$ nor a Dini modulus for
$\sigma$ is required.
\end{theorem}
Compactness makes the positivity uniform. This hypothesis is additional:
we do not prove that an arbitrary continuous right-hand side has a solution
with such a Hessian. The derivative conclusion is G\^ateaux differentiability,
without a claim of Fr\'echet differentiability.

\section{Linearization of the actual bounded equation}
\begin{lemma}\label{lem:actual}
Under the hypothesis of Theorem~\ref{thm:criterion},
$u\in W^{2,p}(X)$ for every finite $p>1$ and has a $C^{1,\alpha}$
representative for every $0<\alpha<1$. Its actual bounded product is
$\sigma^n$. For each smooth real $h$, sufficiently small signed $t$
give bounded genuine admissible $u+th$, and every $\nu\in\cA_0(F)$ satisfies
\begin{equation}\label{eq:stationarity}
 \int_X L_\sigma h\,\mathrm d\nu=0,
 \qquad L_\sigma h=\tr_\sigma(\ddc h).
\end{equation}
\end{lemma}
\begin{proof}
In a coordinate ball the Euclidean Laplacian of $u$ is a constant multiple
of the continuous trace of $\ddc u$. Local constant-coefficient elliptic
regularity gives $W^{2,p}_{\mathrm{loc}}$ for every finite $p$.
One can subtract the Newtonian potential of a cut-off right-hand side;
its high-frequency second derivatives are order-zero multipliers, and
its low-frequency part has locally bounded second-derivative kernels.
The remaining distribution is harmonic. Sobolev embedding and a finite
cover give the asserted regularity. Its continuous representative agrees
with the genuine quasi-psh representative by canonical uniqueness.

Choose smooth $\rho$ with $\ddc\rho\geq\omega$ on a larger coordinate
ball. The genuine bounded psh function $u+\rho$ has decreasing radial
regularizations \cite{Dem}. Their complex Hessians converge uniformly
on smaller balls since $\ddc(u+\rho)$ is continuous. Bedford--Taylor
continuity \cite{BT} and the finite smooth correction therefore identify
the actual product with $\sigma^n$. The same argument works for $u+th$;
uniform positivity makes it admissible for small $|t|$.

Define the actual continuous logarithmic density
\[
 F_t=F+\log\frac{(\sigma+t\ddc h)^n}{\sigma^n}.
\]
The potential $u+th$ solves the equation at $F_t$ with the same scalar
$c_0(F)$. Uniqueness of the scalar gives $B_0(F_t)=B_0(F)$. Thus
\eqref{eq:support} implies $\int(F_t-F)\,\mathrm d\nu\leq0$.
Uniform ellipticity gives the uniform expansion
$F_t-F=tL_\sigma h+O(t^2)$; using both signs of $t$ proves
\eqref{eq:stationarity}. No multiplication of a rough current by an
unspecified weight occurs. In holomorphic coordinates $L_\sigma$ has
a continuous uniformly positive real principal matrix and no lower-order
terms.
\end{proof}

\section{The continuous-coefficient Fredholm argument}
The next lemma supplies the extension from smooth elliptic theory. Its
conclusion is a proof obligation, rather than an assumed rough adjoint
uniqueness theorem.

\begin{lemma}\label{lem:fredholm}
For a continuous positive Hermitian form $\sigma$ and $p>2n$,
\[
 L_\sigma:W^{2,p}(X)\longrightarrow L^p(X,\dV)
\]
is Fredholm of index zero, with kernel the constants. Its range has a
one-dimensional annihilator in $L^{p'}(\dV)$, where $p'=p/(p-1)$.
\end{lemma}
\begin{proof}
Write $m=2n$. Fix a smooth positive Hermitian form $h$ and consider
$L_s=(1-s)L_h+sL_\sigma$, $0\leq s\leq1$. All these operators have
uniform ellipticity and a common modulus of continuity for their principal
coefficients on a finite holomorphic cover.

\emph{Local estimates.}
For a constant positive matrix $a_0$, the smooth multipliers
\[
 \frac{\xi_i\xi_j}{1+a_0^{kl}\xi_k\xi_l}
\]
satisfy the order-zero derivative bounds, including at $\xi=0$.
The multiplier theorem \cite[Chapter 13, Theorem 5.1]{Taylor} gives
\[
 \|D^2v\|_p\leq C(\|a_0^{ij}D_{ij}v\|_p+\|v\|_p)
 \quad(v\in C_c^\infty).
\]
Freeze the actual continuous principal matrix at a point and take the
chart small enough to absorb its oscillation times $\|D^2v\|_p$.
Cut-off differentiation adds only first and zeroth derivatives. The
elementary mollifier interpolation estimate
\[
 \|Dv\|_p\leq\varepsilon\|D^2v\|_p+C_\varepsilon\|v\|_p
\]
absorbs the first derivatives: compare $Dv$ with its mollification at
scale proportional to $\varepsilon$, and bound the derivative of the
mollified function by that scale's inverse times $\|v\|_p$.
Summing the localized estimates and using smooth density proves
\begin{equation}\label{eq:estimate}
 \|v\|_{W^{2,p}}\leq C(\|L_sv\|_p+\|v\|_p),
\end{equation}
with a constant uniform in $s$.

\emph{Maximum principles at this regularity.}
The smooth Alexandrov--Bakelman--Pucci bound on a bounded annulus, for
the fixed continuous coefficient matrix, has the consequence
\[
 \min_{\partial\Omega}w\geq0,\quad L_sw\leq0
       \quad\Longrightarrow\quad w\geq0
\]
for $w\in W^{2,p}(\Omega)$. Indeed smooth approximation in $W^{2,p}$
converges uniformly up to the smooth boundary; the positive parts of the
$L_s$ errors converge in $L^m$, since $p>m$. Apply the smooth ABP bound
to the negative of the approximant, including its vanishing boundary
error, and pass to the limit. The printed ABP proof
\cite[Chapter 14, Proposition 13.1 and Lemma 13.2]{Taylor} uses no
derivative of the coefficients.

If a $C^1$, $W^{2,p}$ function $v$ with $L_sv\geq0$ has an interior
maximum $M$ on a connected open set, the usual tangent-ball argument
proves the strong maximum principle at exactly this regularity.
For completeness, if $w=M-v$ is positive somewhere and zero somewhere,
choose a ball $B_r(x_c)\subset\{w>0\}$, compactly inside a coordinate
chart, tangent to $\{w=0\}$. On its annulus the barrier
\[
 \phi(x)=e^{-k|x-x_c|^2}-e^{-kr^2}
\]
has $L_s\phi>0$ for large $k$, by ellipticity. Choose $\epsilon>0$
so that $w\geq\epsilon\phi$ on the inner sphere; this also holds on
the outer sphere. Weak minimum comparison yields it throughout the
annulus. At the tangent point it contradicts $Dw=0$ by forcing a strictly
positive inward directional derivative. A nontrivial maximum set would
have a boundary point in the connected open set, so such a tangent ball
would exist. Hence $v$ is constant. Applied on $X$, this proves
$\ker L_s=\mathbb R$. If $(L_s-1)v=0$ and $v$ has a positive maximum,
apply this argument on the component of $\{v>0\}$ containing that
maximum. Constancy there contradicts $L_sv=v>0$. Apply it to $-v$
to exclude a negative minimum. Thus $\ker(L_s-1)=0$.

\emph{Invertibility and index.}
Rellich compactness \cite[Chapter 13, equation (6.9)]{Taylor},
\eqref{eq:estimate} and the trivial kernel imply the uniform estimate
\begin{equation}\label{eq:invertestimate}
 \|v\|_{W^{2,p}}\leq C\|(L_s-1)v\|_p.
\end{equation}
Otherwise normalize a contradicting sequence in $W^{2,p}$, take
$s_j\to s$ and a strongly $L^p$ convergent subsequence, and use
\eqref{eq:estimate} and operator-norm continuity in $s$ to get strong
$W^{2,p}$ convergence. The limit is a nonzero element of
$\ker(L_s-1)$, a contradiction.

The smooth starting operator $L_h-1$ is an isomorphism by classical
smooth elliptic index theory and the maximum principle. Its index is
zero by deformation to a smooth Laplacian with a negative zeroth-order
term. The method of continuity gives invertibility for every $s$:
openness is the Neumann lemma; closedness follows from
\eqref{eq:invertestimate}, weak $W^{2,p}$ compactness and strong $L^p$
compactness of solutions for each fixed right-hand side. The estimate
then upgrades convergence to $W^{2,p}$.
Finally $L_\sigma$ differs from the isomorphism $L_\sigma-1$ by the
compact inclusion $W^{2,p}\hookrightarrow L^p$. Compact perturbation
preserves the Fredholm index. Since its kernel has dimension one,
its range has codimension one; $L^p$ duality proves the stated
annihilator conclusion.
\end{proof}

\section{Uniqueness and the derivative}
\begin{proof}[Proof of Theorem~\ref{thm:criterion}]
Choose
\[
 q=1+\frac1{4n},\qquad p=\frac q{q-1}=4n+1>2n.
\]
The imported bound applies since $1<q<n/(n-1)$. Every support has an
$L^q$ density. Lemma~\ref{lem:actual}, smooth density in $W^{2,p}$ and
H\"older duality extend its stationarity to every $W^{2,p}$ test.
Lemma~\ref{lem:fredholm} puts all those densities in a one-dimensional
space. Their integrals equal one, so they coincide. Nonemptiness makes
this unique normalized density an actual support for every continuous
$G$ in \eqref{eq:support}.

For completeness, the derivative argument is the usual support argument.
Monotonicity and the constant-shift law make $B_0$ 1-Lipschitz in the
uniform norm. For continuous $H$, take $\nu_t\in\cA_0(F+tH)$.
Every weak cluster point as $t\to0$ is in $\cA_0(F)$: pass to the limit
in the supporting inequality for each fixed continuous $G$.
Compactness of probabilities and uniqueness imply $\nu_t\to\nu_F$
weakly. For $t>0$, the two supporting inequalities give
\[
 -\int H\,\mathrm d\nu_t
 \leq\frac{B_0(F+tH)-B_0(F)}t
 \leq-\int H\,\mathrm d\nu_F.
\]
The inequalities reverse appropriately for $t<0$. Taking limits proves
the formula for every $H$.
\end{proof}

\section{Explicit products with an arbitrary Hermitian factor}
\begin{proposition}\label{prop:product}
Let $(Y,\kappa)$ be any connected compact Hermitian $k$-fold, $k\geq1$,
normalized by $\int_Y\kappa^k=1$. On $X=\mathbb P^1\times Y$ let
$w$ be Fubini--Study area of mass one, $a>0$, $b>0$ smooth on
$\mathbb P^1$, and $A>0$ continuous there with $\int A w=a$.
Set
\[
 \omega=aw+b\kappa,\qquad \dV=w\wedge\kappa^k,
 \qquad F=\log((k+1)Ab^k).
\]
Then $c_0(F)=1$ and
\begin{equation}\label{eq:productnu}
 \cA_0(F)=\left\{\frac A a w\otimes\nu_\kappa\right\},
\end{equation}
where $\nu_\kappa$ is the smooth positive stationary probability of
$L_\kappa=\tr_\kappa\ddc$. For $k=1$, $\nu_\kappa=\kappa$.
\end{proposition}
\begin{proof}
The zero-mean Green solution on $\mathbb P^1$ of
$\ddc u=(A-a)w$ is bounded and genuine, with continuous complex Hessian.
Its pullback gives the positive continuous form
$\sigma=Aw+b\kappa$ and the actual equation
\[
 (\omega+\ddc u)^{k+1}=(k+1)Ab^k w\wedge\kappa^k=e^F\dV.
\]
Scalar uniqueness gives $c_0(F)=1$. For $k\geq2$, smooth
Gauduchon theory \cite[Theorem 2]{Gaud} gives a smooth positive
stationary probability $\nu_\kappa$; for $k=1$ this is just Stokes.
The block diagonal form gives
$L_\sigma=A^{-1}L_w+b^{-1}L_\kappa$.
For every smooth function depending on all variables, integration against
the probability in \eqref{eq:productnu} annihilates the first term by
base Stokes and the second by vertical stationarity and Fubini.
This candidate has a bounded density, hence lies in the annihilator of
Lemma~\ref{lem:fredholm}. Theorem~\ref{thm:criterion} supplies an actual
support in that one-dimensional space. Normalization identifies the two,
so the candidate satisfies the inequality for every continuous $G$.
Testing only separated functions is not used to establish support.
\end{proof}

\subsection{A non-Dini density on a non-K\"ahler threefold}
Take $Y=(\mathbb C^2\setminus\{0\})/\langle z\mapsto2z\rangle$ and
$\kappa$ a unit-volume multiple of $|z|^{-2}\ddc|z|^2$.
It descends to the quotient and is positive. Writing $s=|z|^2$,
the identity $s(s^{-1})''+2(s^{-1})'=0$ gives $\ddc\kappa=0$.
Consequently $\nu_\kappa=\kappa^2$, since
$L_\kappa h\,\kappa^2=2\ddc h\wedge\kappa$ has integral zero.
The smooth manifold $Y$ is $S^3\times S^1$ and has $H^2(Y,\mathbb R)=0$.
A closed positive $(1,1)$-form would be exact but have positive square
integral, contradicting Stokes. Thus $Y$ and $X=\mathbb P^1\times Y$
are non-K\"ahler, the latter by restriction to the factor.

Let $\tau\in[-1,1]$ be the usual height on $\mathbb P^1$, put
$b=2+\tau$, and define
\[
 f(t)=\begin{cases}1/\log(e/|t|),&0<|t|\leq1,\\0,&t=0,\end{cases}
 \qquad A=a+\epsilon\left(f(\tau)-\int f(\tau)w\right),
 \quad0<\epsilon<a.
\]
Then $A$ is continuous, uniformly positive and has mean $a$.
Near the equator its modulus is bounded below by a positive multiple
of $1/\log(e/r)$, whose Dini integral diverges. It is not H\"older of
any positive order. The same is true of $F=\log(3Ab^2)$: the smooth
factor contributes only $O(r)$, and cannot cancel that modulus.
The background $\omega=aw+b\kappa$ is genuinely nonclosed because
the component $db\wedge\kappa$ is nonzero. Proposition~\ref{prop:product}
gives the explicit supporting probability $(A/a)w\wedge\kappa^2$.
The reference is zero and $e^F$ is strictly positive everywhere.

\section{Proof boundary}
The imported foundations are canonical bounded psh theory and bounded
Bedford--Taylor continuity; constant-coefficient multiplier estimates;
Sobolev embedding, smooth density and Rellich compactness; the smooth
local ABP bound; smooth elliptic regularity and Fredholm theory; and
standard Banach-space continuity and compact-perturbation index arguments.
The continuous-coefficient Fredholm extension is proved in
Lemma~\ref{lem:fredholm}, rather than quoted as a new imported premise.
Gauduchon's theorem is applied only to the smooth factor $\kappa$.
No nonsmooth-coefficient Schauder theorem, rough current product,
derivative of the principal coefficients or Markov-generator core is
assumed.

The original P16 asks for uniqueness at every continuous $F$ on every
compact Hermitian manifold. This article has an additional continuous
positive Hessian hypothesis and supplies no exceptional density with
two supports. Therefore it does not close original P16.

\begin{thebibliography}{9}
\bibitem{BT} E. Bedford and B. A. Taylor,
\emph{A new capacity for plurisubharmonic functions},
Acta Math. 149 (1982), 1--40, Theorem 2.1 and its printed proof.
\bibitem{Dem} J.-P. Demailly, \emph{Complex analytic and differential geometry},
Chapter I, canonical psh representatives and radial regularization,
\url{https://www-fourier.univ-grenoble-alpes.fr/~demailly/manuscripts/agbook.pdf}.
\bibitem{Gaud} P. Gauduchon,
\emph{Les m\'etriques standard d'une surface complexe compacte \`a
premier nombre de Betti pair}, Ast\'erisque 126 (1985), 129--135,
Theorem 2 (the general Hermitian theorem, independent of the parity
assumption in Theorem 1),
\url{https://www.numdam.org/item/AST_1985__126__129_0/}.
\bibitem{Taylor} M. E. Taylor, \emph{Partial differential equations III},
author-hosted Chapters 13 and 14: order-zero multipliers, compact Sobolev
inclusions, and the ABP estimate and contact-set proof.
\bibitem{S07} \emph{Singular non-Gauduchon product adjoints},
checked programme article,
\url{https://github.com/MingchenXia/non-kahler-pluripotential-theory/blob/main/research/papers/S07-singular-non-Gauduchon-product-adjoints.tex}.
\end{thebibliography}
\end{document}
